\begin{definition}\label{def:ParityVector}\lean{CollatzMapBasics.ParityVector}\leanok
A parity vector is a finite sequence of bits (0 and 1). In the formalization, it is represented as a list of boolean values, where \texttt{false} corresponds to 0 and \texttt{true} corresponds to 1.
\end{definition}

\begin{definition}\label{def:ParityVector_q}\lean{CollatzMapBasics.q}\leanok
For a parity vector $v$, $q(v)$ is the number of ones (true entries) in $v$.
\end{definition}

\begin{definition}\label{def:ParityVector_size}\lean{CollatzMapBasics.size}\leanok
The size of a parity vector $v$, denoted $|v|$, is its length.
\end{definition}

\begin{definition}\label{def:ParityVector_onesRatio}\lean{CollatzMapBasics.onesRatio}\leanok
The ones-ratio of a parity vector $v$ is the rational number $q(v) / |v|$. If the vector is empty, the ratio is defined to be 0.
\end{definition}

\begin{definition}\label{def:ParityVector_V}\lean{CollatzMapBasics.V}\leanok
For $j, n \in \mathbb{N}$, the parity vector $V(j, n)$ of length $j$ for the compact Collatz
sequence starting at $n$ is defined as
\[
    V(j, n) = (X(n), X(T(n)), \dots, X(T^{j-1}(n))).
\]
\end{definition}

\begin{lemma}\label{lemma:V_length}\lean{CollatzMapBasics.V_length}\leanok
The length of $V(j, n)$ is $j$.
\end{lemma}
\begin{proof}\leanok
Immediate from the definition.
\end{proof}

\begin{lemma}\label{lemma:V_get}\lean{CollatzMapBasics.V_get}\leanok
For any $i < j$, the $i$-th entry of $V(j, n)$ is $X(T^i(n))$.
\end{lemma}
\begin{proof}\leanok
Immediate from the definition.
\end{proof}

\begin{definition}\label{def:ParityVector_Precedes}\lean{CollatzMapBasics.Precedes}\leanok
The elementary precedence relation $\prec_e$ is defined by swapping a $01$ subword with $10$:
\begin{equation}
    w_1 01 w_2 \prec_e w_1 10 w_2.
\end{equation}
The precedence relation $\preceq$ is the reflexive transitive closure of $\prec_e$.
\end{definition}

\begin{definition}\label{def:E_vec}\lean{CollatzMapBasics.E_vec}\leanok
For $k, n \in \mathbb{N}$, the vector $E(k, n)$ is the function from $\{0, \dots, k-1\}$ to $\{0, 1\}$ defined by
\begin{equation}
    E(k, n)(i) = X(T^i(n)).
\end{equation}
\end{definition}

\begin{lemma}\label{lemma:E_vec_le_one}\lean{CollatzMapBasics.E_vec_le_one}\leanok
Each entry of $E(k, n)$ is at most 1.
\end{lemma}
\begin{proof}\leanok
Immediate from the definition of $X(n)$.
\end{proof}

\begin{lemma}\label{lemma:num_odd_steps_eq_E_vec_sum}\lean{CollatzMapBasics.num_odd_steps_eq_E_vec_sum}\leanok
The number of odd steps $Q(k, n)$ is equal to the sum of the entries of $E(k, n)$:
\begin{equation}
    Q(k, n) = \sum_{i=0}^{k-1} E(k, n)(i).
\end{equation}
\end{lemma}
\begin{proof}\leanok
By definition of $Q(k, n)$ and $E(k, n)$.
\end{proof}

\begin{lemma}\label{lemma:E_vec_restrict}\lean{CollatzMapBasics.E_vec_restrict}\leanok
If $E(k+1, m) = E(k+1, n)$, then $E(k, m) = E(k, n)$.
\end{lemma}
\begin{proof}\leanok
The agreement on the first $k+1$ entries implies agreement on the first $k$ entries.
\end{proof}

\begin{lemma}\label{lemma:num_odd_steps_eq_of_E_vec_eq}\lean{CollatzMapBasics.num_odd_steps_eq_of_E_vec_eq}\leanok
If $E(k, m) = E(k, n)$, then $Q(k, m) = Q(k, n)$.
\end{lemma}
\begin{proof}\leanok
Directly from Lemma~\ref{lemma:num_odd_steps_eq_E_vec_sum}.
\end{proof}
